\section{Conclusion}
\label{sec:conclusion}

This study examines whether document-level sentiment extracted from Japanese
MD\&A disclosures is associated with abnormal stock returns around EDINET
filing dates, and whether GPT-based sentiment contains incremental explanatory
power beyond a finance-specific Japanese lexicon. Using 1,087 annual filings
from Nikkei 225 firms over 2018--2023, we compare GPT-based document sentiment
with a Japanese adaptation of the Loughran--McDonald Master Dictionary in a
cross-sectional event-study framework with year and industry fixed effects and
firm-clustered standard errors.

The evidence shows that disclosure tone is associated with abnormal returns, but that this relation is captured more consistently by the
Japanese LMMD measure than by the GPT-based document score. GPT sentiment is
positively associated with CARs in some shorter-window standalone
specifications, but its coefficient attenuates once finance-specific lexical
tone is included. By contrast, LMMD net tone remains economically meaningful
and statistically more stable in joint specifications, especially in the
primary $[-1,1]$ event window. Wider-window results reinforce this pattern but
are interpreted as robustness evidence because longer windows are more exposed
to unrelated firm news.

When cumulative abnormal returns are used as the benchmark for economically relevant disclosure tone, the Japanese LMMD measure exhibits the stronger and more persistent association in this setting, while GPT contributes little incremental return-relevant variation once LMMD tone is included. Methodologically, the paper contributes a
validated Japanese LMMD construction, a large-scale GPT-based scoring pipeline
for Japanese MD\&A text, and a unified event-study design that allows direct
comparison between dictionary-based and LLM-based sentiment measures.

The results have implications for applied NLP in finance. For standardized
annual disclosures, transparent domain dictionaries may capture much of the
return-relevant sentiment signal, while LLM-based measures may add more value
in settings where language is less standardized, more interactive, or more
context-dependent, such as earnings calls, press releases, management guidance,
or richer narrative sections of filings. Future work could examine section-level
sentiment, alternative prompting and aggregation strategies, higher-frequency
disclosure events, and cross-country comparisons. Several limitations frame the
interpretation: the analysis focuses on annual MD\&A filings, market reactions,
and document-level sentiment in the Japanese institutional context; finer-grained
textual modeling and alternative return benchmarks may reveal additional
channels linking disclosure language to prices.